\documentclass[11pt,a4paper]{article}
\usepackage[T1]{fontenc}
\usepackage{lmodern,amsmath,amssymb,amsthm,mathtools,booktabs,microtype}
\usepackage[margin=27mm]{geometry}
\usepackage[colorlinks=true,linkcolor=blue,citecolor=blue,urlcolor=blue]{hyperref}
\newtheorem{theorem}{Theorem}[section]
\newtheorem{lemma}[theorem]{Lemma}
\newtheorem{proposition}[theorem]{Proposition}
\newtheorem{corollary}[theorem]{Corollary}
\newcommand{\Q}{\mathbb Q}\newcommand{\Z}{\mathbb Z}
\DeclareMathOperator{\Tr}{Tr}\DeclareMathOperator{\rad}{rad}
\setlength{\emergencystretch}{2em}
\title{A Thirteen-Variable Family under the Rational $abc$ Conjecture}
\author{}\date{9 November 2026}
\begin{document}\maketitle
\begin{abstract}
Assuming the rational $abc$ conjecture, together with the specified
elliptic group, pole-realization and simultaneous prime-pattern inputs,
we construct a family of integer polynomials in thirteen total variables.
Every sufficiently large fixed odd input shift gives rational input
projection exactly $\Z$. Each member has total degree less than
$2^{44}$. A logarithmic height estimate for odd contact radicals and
a height estimate for a sum of sixth powers allow the rational input
and one elliptic index to be bounded together. Only three elliptic
points and two pole certificates occur. All six local tests have order
one and use the elliptic formal parameter itself. The coefficient
threshold is ineffective. The additional $abc$ hypothesis is essential
to the proof given here; this is not a reduction under the arithmetic
premises of the separate twenty-seven-variable theorem alone.
\end{abstract}

\section{Statement and hypotheses}

All variable counts include the input. For a rational number $a/b$
in lowest terms with $b>0$, write
$h(a/b)=\log\max(|a|,b)$.
The rational $abc$ conjecture means that for every $\epsilon>0$
there exists $K_\epsilon\ge1$ such that
\[
 c\le K_\epsilon\rad(abc)^{1+\epsilon}
 \quad(a+b=c,\ a,b,c\text{ positive and pairwise coprime}).
\]

Set
\[
 K=\Q(\sqrt2),\quad D=(3-\sqrt2)^2,\quad
 E:y^2=x^3-Dx,
\]
\[
 P=\left(\frac{57}{8}-\frac72\sqrt2,
                 -\frac{165}{8}+\frac{417}{32}\sqrt2\right).
\]
In addition to $abc$, retain the following arithmetic inputs with
exactly their stated scope in~\cite{Dependencies,Source27}.
\begin{enumerate}
\item $2E(K)=\Z P$, with the indicated non-torsion generator
and the source index-height bound. The direct point chart below
therefore describes the nonzero multiples of $P$.
\item The pole system, for $b\ne0$ and $B_0=17b$,
\[
 p^2=B_0(1-e)(1+3e),\qquad
 B_0q^2=(s^2-B_0e(B_0-b^2))(s^2+3B_0e(B_0-b^2)),
\]
excludes negative odd valuations of $b$ outside $\{2,3,17\}$.
For a prepared input
\begin{equation}\label{eq:prepared}
 \begin{gathered}
 0<b<17,\quad b\text{ has square reduced denominator},\quad
 b\in\Q_2^{\times2}\cap\Q_3^{\times2}\cap\Q_{17}^{\times2},\\
 v_2(b)\ge3,\quad v_3(b)\ge1,\quad v_{17}(b)\ge2,
 \end{gathered}
\end{equation}
the retained restricted-converse and descent hypotheses supply a
witness with $pqs\ne0$. The nonzero-witness consequence and its
single conic chart are proved in~\cite[Lemma~2]{Source27}.
\item The simultaneous prime-pattern input supplies the specified
prime contact numerators, Gaussian congruences, conjugate reduction
orders, fixed divisibility and avoidance, and compatible real
approximation for any fixed finite set of positive rational square
scales. Its precise quantified version and fixed-avoidance deduction
from Sun's theorem are given in~\cite{Dependencies}. The order of
choices is respected in Section~\ref{sec:complete}.
\end{enumerate}
The exponent-nineteen contact-height input of the earlier source is
not needed here. It is replaced by Lemma~\ref{lem:contact}, which
uses the additional $abc$ assumption.

The proved one-point case of the aggregate trace-height theorem
\cite{Trace} gives constants $A>0$, $J\ge0$, depending only on
$E,P$, such that
\begin{equation}\label{eq:traceheight}
 n^2\le A h(T(n))+J,\qquad
 T(n)=\Tr_{K/\Q}x([n]P)\quad(n\in\Z\setminus\{0\}).
\end{equation}
This applies von Buhren's theorem after checking the absence of
constant abelian translates for the fixed trace function. Its additive
constant is ineffective. It is a proved height theorem, not an
additional conjecture. For a nonzero multiple of $P$, the abscissa
is a positive square in both real embeddings; the curve equation
then gives $x>3\mp\sqrt2$ in the respective embeddings, hence
$T(n)>6$.

\begin{theorem}\label{thm:main}
Under the assumptions above, there is a computable family
\[
 F_B\in\Z[X,Y_1,\ldots,Y_{12}],\qquad B\ge3\text{ odd},
\]
with the following properties.
\begin{enumerate}
\item Every member is given by a complete finite integer coefficient
specification and satisfies
\[
             \deg F_B\le13155996618240<2^{44}.
\]
\item There exists a constant $B_*$ such that every odd $B\ge B_*$
has rational input projection exactly $\Z$.
\item No numerical sufficient value of $B_*$ is determined.
The supplied member $B=3$ is a coefficient-verification fixture,
not a polynomial certified to define $\Z$.
\end{enumerate}
The integer $B$ is a fixed coefficient parameter and is never a
quantified coordinate.
\end{theorem}

\section{Two consequences of $abc$}

Put
\[
 \Theta=\left\{0,-\frac12,-1,\frac1{2399},-\frac{32}{407}\right\},
 \qquad S_0=\{p:p\le76176\}.
\]
For $t\in\Q\setminus\Theta$, let $M_0(t)$ be the product of
primes outside $S_0$ for which $v_p(t-\theta)$ is positive
and odd for some $\theta\in\Theta$.

\begin{lemma}\label{lem:contact}
Assume rational $abc$. There is a constant $C_0\ge0$ such that
\[
                    h(t)\le\frac{10}{9}\log M_0(t)+C_0
                  \qquad(t\in\Q\setminus\Theta).
\]
\end{lemma}
\begin{proof}
For a primitive pair $(u,v)$, consider the squarefree homogeneous
binary form
\[
 G(u,v)=vu(u-v)(u-2401v)(32u+343v)
\]
of degree five. Granville's consequence of $abc$
\cite[Proposition~1]{Granville}, with $\epsilon=1/20$, gives
\[
 \log\rad(G(u,v))\ge\frac{59}{20}\log\max(|u|,|v|)-C.
\]
For a nonzero integer $N$, if $\operatorname{sf}(N)$ denotes
the product of primes having odd valuation in $N$, then
$\rad(N)^2\le |N|\operatorname{sf}(N)$, prime by prime.
Outside $S_0$, two of the five primitive factors of $G$ cannot
vanish simultaneously. Thus its squarefree part outside $S_0$
is precisely the odd contact radical for
$\infty,0,1,2401,-343/32$. The omitted primes contribute a fixed
bounded factor. Also $|G(u,v)|\le C'\max(|u|,|v|)^5$.
Combining the two estimates gives
\[
 \log M_5(u/v)\ge\frac9{10}h(u/v)-C''.
\]
The integral linear fractional transformation $t\mapsto2+1/t$
has determinant $-1$, sends the contacts $\Theta$ to the five
original contacts, and changes logarithmic height by at most
$\log3$. Its contact radicals coincide outside $S_0$.
These assertions follow directly by substituting a primitive
numerator-denominator pair; their fixed exceptional determinants
have prime divisors in $S_0$. Absorbing the fixed constants proves
the claim.
\end{proof}

\begin{lemma}\label{lem:sixth}
Assume rational $abc$. There is a constant $C_1\ge0$ such that
\[
 h(u)+h(v)\le\frac45 h(u^6+v^6)+C_1\qquad(u,v\in\Q).
\]
One may take
$C_1=\frac7{15}\log2+\frac{11}{15}\log K_{1/11}$.
For $(u,v)\ne(0,0)$ and $a=(u^6+v^6)^2$, equivalently,
\[
                       h(u)+h(v)\le\frac25h(a)+C_1.
\]
\end{lemma}
\begin{proof}
The assertion is immediate if a coordinate is zero, since
$h(w^6)=6h(w)$. Otherwise choose integers $A',B',D_0$, with
$D_0>0$ and $\gcd(D_0,A',B')=1$, representing
$u=A'/D_0$, $v=B'/D_0$. Write $g=\gcd(|A'|,|B'|)$ and
$A_0=A'/g$, $B_0=B'/g$. Then $\gcd(A_0,B_0)=\gcd(g,D_0)=1$.
Put $H=\max(|A_0|,|B_0|)$ and $N=A_0^6+B_0^6$, so
$H^6\le N\le2H^6$.

Let $k=\gcd(N,D_0^6)$, $m_0=N/k$, and $n=D_0^6/k$.
The fraction $s=u^6+v^6=m/n$, where $m=g^6m_0$, is reduced.
Put $M=\max(m,n)$, so $h(s)=\log M$. Every prime dividing
$N$ either divides $D_0$ or survives in $m_0$, giving
$\rad(N)\le D_0m_0$. Apply $abc$ at $\epsilon=1/11$ to the
pairwise coprime triple $A_0^6,B_0^6,N$:
\[
 H^{7/2}\le K_{1/11}^{11/12}D_0m_0.
\]
Since $D_0^6=nN/m_0\le2nH^6/m_0$, this implies
\[
 H^{5/2}\le2^{1/6}K_{1/11}^{11/12}m_0^{5/6}n^{1/6}.
\]
Moreover
\[
 \max(gH,D_0)\le2^{1/6}H(M/m_0)^{1/6}.
\]
Taking logarithms and combining these inequalities gives
\begin{align*}
 \log\max(gH,D_0)
 &\le\frac7{30}\log2+\frac{11}{30}\log K_{1/11}
       +\frac16\log m_0+\frac1{15}\log n+\frac16\log M\\
 &\le\frac25\log M+\frac12 C_1.
\end{align*}
Each of $h(u),h(v)$ is at most the left-hand side. Adding proves
the first assertion; $h(s^2)=2h(s)$ proves the second.
The common numerator factor and all denominator cancellations
are thus included.
\end{proof}

\section{The complete rational source}

Fix an odd integer $B\ge3$ and set
\begin{equation}\label{eq:input}
                            \beta=2X^2+B.
\end{equation}
For rational $X=u/v$ in lowest terms, integrality of $\beta$
implies $v^2\mid2$, hence $v=1$. Thus $\beta\in\Z$ if and
only if $X\in\Z$. For integer input $\beta$ is odd, and
$\beta\ge B$ for every rational input.

Retain three pairs of rational coordinates, one for each of the
points $A,H,U$. For any such pair $(h_0,h_1)$ put
\[
 h=h_0+h_1\sqrt2,\quad x=h^2,\quad j^2=h^4-D,\quad y=hj.
\]
The direct point-chart theorem of~\cite{Source27} proves that
these are exactly the nonzero multiples of $P$. Each finite
choice $j$ is encoded by one scalar monic quartic as follows.
Write $h^4-D=a+b\sqrt2$, put
$c=2a+b$, $d=-a-b$, and write
$(c+d\sqrt2)^2(a+b\sqrt2)=a'+b'\sqrt2$.
Then
\[
 Y^4-4a'Y^2+8(b')^2=0,\quad
 j=\frac{Y/2+(b'/Y)\sqrt2}{c+d\sqrt2}.
\]
Its discriminant is
$2^{15}(b')^2(2a^2-b^2)^4(a^2-2b^2)^2$, nonzero at every
rational pair $(h_0,h_1)$. Indeed $h^4-D\ne0$ by its norm,
the two quadratic norm factors are anisotropic over $\Q$, and
$b'=-4a^3+6ab^2+3b^3$ is anisotropic by Eisenstein at 3.
The square-root chart is complete because its transformed rational
component is $2(u^3-2v^3)$ for $j=u+v\sqrt2$, which is nonzero
for $j\ne0$. These are algebraic facts independent of $abc$.

Compute $V_0=H-P$ and $V_1=U-A$ by the usual affine addition
formulas. Put $f=1/x$, $z=-x/y$, let $\sigma$ be conjugation
of $K/\Q$, and define
\[
 L=\frac{z(U)}{z(H)},\qquad
 L'=\sigma\!\left(\frac{z(V_1)}{z(V_0)}\right).
\]
Use precisely the six field-valued tests
\begin{equation}\label{eq:tests}
 D_1=f(H),\ D_2=f(U),\ D_3=\sigma f(V_0),\
 D_4=\sigma f(V_1),\ D_5=L-L',\ D_6=L-\beta.
\end{equation}
Every test is assigned order one. No scalar-height pole or Lucas
relation is present.

Let $T=\Tr x(A)$, and use the two positive rational square scales
\[
                         a_0=1,\qquad a_1=(\beta^6+T^6)^2.
\]
Retain one rational slope $\lambda$. For $r=0,1$, set
\begin{align}
 t_r&=a_r\lambda,&\delta_r&=t_r-1,&
 F_r&=\prod_{\theta\in\Theta}(t_r-\theta),&
 R_r&=\frac{F_r}{a_r^5\delta_r^6},\label{eq:detector}\\
 W_r&=\prod_{j=1}^6\left(1+
             \frac{D_j^2}{\delta_r^4R_r}\right),&
 b_r&=\frac{R_r}{\Tr(W_r^{-1})^2}.\label{eq:W}
\end{align}
All displayed divisions, affine-chart denominators, $F_r$,
$\delta_r$, $\operatorname{Norm}(W_r)$, and $\Tr(W_r^{-1})$
are guarded nonzero. These conditions are part of the source.

For each $b=b_r$, retain independent rational parameters $\tau_r,s_r$.
Put $B_0=17b$ and
\[
 e=\frac{\tau_r^2-B_0}{\tau_r^2+3B_0},\qquad
 q^2=\frac{(s_r^2-B_0e(B_0-b^2))
                  (s_r^2+3B_0e(B_0-b^2))}{B_0}.
\]
Guard the denominator and require $q\ne0$. This reconstructs
the first pole ordinate as
$p=-4B_0\tau_r/(\tau_r^2+3B_0)$.
Each rational-square equation $q^2=N_0/D_0$ is made monic
by $Z=D_0q$, giving $Z^2=N_0D_0$ with $Z\ne0$ and $D_0\ne0$.
No new unrestricted rational parameter is used for these roots.

\section{The local implication}

\begin{lemma}\label{lem:local}
For any rational solution of the guarded source, write
$A=[n]P$, $H=[h]P$, $U=[u]P$.
If $\beta\ne n$ and $N$ is the positive absolute numerator
of $\beta-n$ in lowest terms, then
\[
                         M_0(t_r)\mid N\qquad(r=0,1).
\]
\end{lemma}
\begin{proof}
Fix a prime $p\mid M_0(t_r)$ and abbreviate $R=R_r$, $W=W_r$.
Outside $S_0$ only one contact factor has positive valuation,
$\delta_r$ is a unit, and $v_p(R)$ is odd because $a_r$ is
a square. For each place $\mathfrak p$ above $p$, the two terms
of a factor of $W$ have valuations of different parity. Hence
\[
 v_{\mathfrak p}(W)=\sum_j
       \min\{0,2v_{\mathfrak p}(D_j)-v_p(R)\}\le0.
\]
Put $w=\max_{\mathfrak p\mid p}v_{\mathfrak p}(W)\le0$.
Then $v_p\Tr(W^{-1})\ge-w$.
The valuation of $b=R/\Tr(W^{-1})^2$ is odd, so the pole
obstruction gives
\[
                   0<v_p(b)\le v_p(R)+2w.
\]
Thus $\nu=v_p(R)>0$ and $w>-\nu/2$. At a place attaining
$w$, each summand in the valuation of $W$ is at least $w$.
It follows that $v_{\mathfrak p}(D_j)>\nu/4>0$, so all
six tests have valuation at least one at that place.

The first four tests put $H,U$ into the formal group at
$\mathfrak p$ and $V_0,V_1$ into it at the conjugate place.
At these good unramified places the formal logarithm satisfies
\[
                         \log_E(Q)=z(Q)(1+O(p^4)).
\]
Indeed its first omitted term is a fixed multiple of $Df(Q)^2$;
$v_p(f(Q))\ge2$ in the formal group. The remaining coefficient
estimates, uniformly for $p>76176$, are those proved in the
formal-logarithm paper~\cite{Logs}. The exact logarithmic ratios
are $u/h$ and $(u-n)/(h-1)$, respectively.

The affine guards imply $h\ne0,1$. At least one of $h,h-1$
is a $p$-adic unit, so the corresponding exact ratio is integral.
Relative precision four makes its truncated ratio integral.
The test $L-L'\equiv0\pmod p$ makes the other truncated ratio
integral as well; the same relative precision then makes the
other exact ratio integral. Relative errors are now absolute
errors of valuation at least four. Consequently
\[
 L\equiv u/h,\qquad L'\equiv(u-n)/(h-1),\qquad
 L\equiv L'\pmod p.
\]
Multiplying by the integral numbers $h,h-1$ and subtracting yields
$L\equiv n\pmod p$. The target test gives $L\equiv\beta\pmod p$.
In particular $\beta$ is $p$-integral, and $p$ divides the
reduced numerator of $\beta-n$. This holds at every consumed
prime and proves the divisibility. The same argument includes
inert primes; if their formal-group requirements are incompatible,
no source solution occurs at that prime.
\end{proof}

\section{Soundness and the uniform input shift}

Suppose a rational source solution has $\beta\ne n$. By
Lemma~\ref{lem:local} and Lemma~\ref{lem:contact},
\[
 h(\lambda),h(a_1\lambda)\le\frac{10}{9}\log N+C_0.
\]
The quotient-height inequality therefore gives
\[
 h(a_1)\le\frac{20}{9}\log N+2C_0.
\]
Lemma~\ref{lem:sixth}, with $u=\beta,v=T(n)$, implies
\begin{equation}\label{eq:combined}
 h(\beta)+h(T(n))\le\frac89\log N+C,
 \qquad C=\frac45C_0+C_1.
\end{equation}
All constants here depend only on the fixed contact set and the
fixed sixth-power map. They are independent of the shift $B$,
the rational input, and every witness.

Writing $\beta=a/b$ in lowest terms, one has
$N=|a-nb|\le e^{h(\beta)}(1+|n|)$. Substitution in
\eqref{eq:combined} gives
\begin{equation}\label{eq:absorb}
 \frac19h(\beta)+h(T(n))
                  \le\frac89\log(1+|n|)+C.
\end{equation}
With \eqref{eq:traceheight}, this implies
\[
 |n|^2\le\frac{8A}{9}\log(1+|n|)+AC+J
       \le\frac{8A}{9}|n|+AC+J.
\]
Choose an integer $M>8A/9+\sqrt{AC+J}$. Then $|n|<M$,
and~\eqref{eq:absorb} gives
\[
                       h(\beta)\le8\log(1+M)+9C=:H_*.
\]
Choose $B_*>e^{H_*}$. For every fixed odd $B\ge B_*$,
\eqref{eq:input} gives $h(\beta)\ge\log\beta\ge\log B>H_*$,
a contradiction. Thus $\beta=n\in\Z$, and the elementary
input-map argument yields $X\in\Z$.
This proves soundness uniformly for every sufficiently large shift.
It does not determine that shift effectively.

\section{Completeness}\label{sec:complete}

Fix an integer input $X$ and an odd coefficient $B\ge3$.
Then $n=\beta$ is an odd integer at least three. Fix
$A=[n]P$ and its two positive square scales before choosing
any prime. They are distinct because $a_1>1$.
Temporarily put $H=[h]P$ and $U=[hn]P$.

Choose a nonexceptional $h_*\in64\Z$, and work in the real
orbit closure of $64\Z P$. Choose a compact real neighborhood
for $H$ around $[h_*]P$ on which the point coordinates, both difference charts,
and the six tests are defined and bounded. Such a neighborhood
exists: the rational functions restricted to the multiplication
curve $H\mapsto(H,[n]H)$ are not identically singular.
The difference charts exclude $h=\pm1$, while every nonzero
multiple of the non-torsion $P$ has nonzero abscissa and ordinate.
Pass to a smaller compact neighborhood to avoid all remaining
zero denominators. Enlarge $L$ with $64\mid L$ and the real
component-group condition so that the progression
$h\equiv h_*\pmod L$
preserves fixed lower valuation bounds for these tests at
$2,3,17$. These are precisely the local-continuity and real-orbit
steps of the source realization, with fewer point roles and no
scalar constraint.

Choose a fixed $D_0$ with $D_0^2\theta/a_r\in\Z$ for all
contacts and both scales, set $v=D_0^2w^2$, and seek
$\lambda=u/v$. Apply the retained prime-pattern theorem to
the distinct forms $u-(\theta/a_r)v$, after repetitions have
been removed. All bad primes and previously fixed moduli are
excluded before this application. The theorem permits arbitrarily
large divisibility of $w$, a sufficiently large prescribed lower
bound on $u/w^2$, and the conjugate reduction-order conditions.

As in~\cite{Dependencies,Baseline}, the annihilating order moduli
$U_0,V_0$ for the two conjugate families satisfy
\[
 \gcd(U_0,V_0)=1,\qquad \gcd(U_0,L)=16,\qquad\gcd(V_0,L)=1.
\]
Thus the congruences
$h\equiv0\pmod{U_0}$, $h\equiv1\pmod{V_0}$,
$h\equiv h_*\pmod L$ are compatible, since $16\mid h_*$.
Choose $h$ in that
progression while retaining the real neighborhood and avoiding
the finitely many excluded multiples. At one place above each
contact prime, $H,U$ are formal; at the conjugate place,
$H-P,U-A$ are formal. Both exact logarithmic ratios are $n$.
The constant approximation has precision four, so all six tests
have their assigned order one, including the target $n=\beta$.

The inherited detector preparation now applies without alteration.
For clarity, its main ingredients are the following. The identity
\[
 R_r=\frac{\prod_{\theta\in\Theta}(u-(\theta/a_r)v)}
                  {v^5(a_r\lambda-1)^6}
\]
has a rational-square denominator since $v$ is a square.
At each chosen contact prime $v_p(R_r)=1$; the successful
conjugate place gives all inverse norm-factor residues equal to
one, and the trace of $W_r^{-1}$ is a unit with residue one or two.
The resulting $b_r$ has square reduced denominator, and the
prescribed congruences give its local squareclasses at $2,3,17$.
The fixed lower bounds on the tests and sufficient divisibility
of $w$ give the depths in~\eqref{eq:prepared}.
On the compact real set, uniformly,
\[
 R_r\sim a_r^{-6}\lambda^{-1},\qquad
 W_r=1+O(\lambda^{-3}),\qquad
 b_r\sim(4a_r^6\lambda)^{-1}.
\]
Thus $0<b_r<17$ for the sufficiently large lower bound on
$\lambda$ fixed before prime selection. All detector guards
are nonzero. These calculations are the source comparison-pole
preparation proof~\cite{Baseline}; they depend on the finite
list of tests and their orders, not on a Lucas or scalar test.

Each prepared input has an independent pole certificate with
$pqs\ne0$ by the retained pole-realization hypotheses.
Since $p\ne0$ excludes $e=1$, the substitution
$\tau=p/(e-1)$ gives the single chart used above and its
nonzero denominator. Thus both nonzero quadratic-root conditions
are realized. The direct point charts cover $A,H,U$.
Every rational input $X\in\Z$ therefore has all required
witnesses, for every fixed odd $B\ge3$.

\section{The polynomial, its bounds, and verification}

The retained coordinates are exactly
\begin{center}
\begin{tabular}{lr}
\toprule
Coordinates & Number\\
\midrule
Input $X$ & 1\\
Three point pairs for $A,H,U$ & 6\\
Comparison slope $\lambda$ & 1\\
Two independent pole pairs $(\tau_r,s_r)$ & 4\\
\midrule
Retained before finite-root elimination & 12\\
\bottomrule
\end{tabular}
\end{center}
There are three monic point quartics and two monic pole quadratics.
The point algebra has rank $4^3=64$, and the full algebra has
rank $4^3 2^2=256$. Its complete guard is $gq_1q_2$ with
$g$ in the point algebra, and each quadratic radicand also belongs
to that algebra. There are no further residual equations.
The separating-element theorem~\cite{Active} therefore replaces
all five roots and all nonzero guards by one rational coordinate,
without changing the input projection. This gives thirteen total
coordinates. Its determinant formulas and its 64-factor product,
applied to the supplied integer coefficient circuit, define $F_B$
completely.

The native generator weights are
\[
                         6,6,6,50912,76641.
\]
The full guard and its point factor have weights $156672$ and
$29119$. In the notation of~\cite[Theorem~4.1]{Active},
the multiplication-element weight is at most $W_1r+W_0$, where
\[
                    W_1=10200912,\qquad W_0=75650993.
\]
Thus
\[
 \deg F_B\le256\bigl(10200912\cdot64^2
                         +(75650993+1)\cdot127\bigr)
             =13155996618240<2^{44}.
\]
The input shift is a named coefficient parameter of weight zero
in the coefficient graph. No operation or branch depends on its
numerical value. Consequently these weight inequalities, and hence
the degree bound, hold symbolically for every fixed $B$.

The directory \nolinkurl{computation/few_variables/conditional13}
contains the rational-source builder, its integer coefficient
specification, and independent source and bound certificates.
The source has 1731 rational-circuit nodes and its integer
coefficient graph has 748 nodes. The independent conversion auditor
checks every affine substitution, the three point quartics and
their automatic nonzero factors, both pole-square clearings,
and all 46 atomic guard factors by exact expression identities.
The separate finite-algebra checker verifies coefficient dependencies,
companion-column inequalities and all 64 degree cases without
importing the builder.

For the algebraic fixture $B=3$ it certifies
\[
 \log_2 H(F_3)\le31315024824505<2^{45}.
\]
This is not a height bound for a member certified to define $\Z$.
The proof provides an ineffective sufficient shift, whereas the
coefficient algorithms construct and bound each specified member.
The exact finite checks establish their recorded algebraic assertions;
they do not verify $abc$ or the inherited arithmetic premises.

\begin{thebibliography}{9}
\bibitem{Granville}
A.~Granville, Rational and integral points on quadratic twists of a
given hyperelliptic curve, \emph{International Mathematics Research
Notices} (2007), article rnm027, Proposition~1.
\url{https://doi.org/10.1093/imrn/rnm027}; author's manuscript,
\url{https://dms.umontreal.ca/~andrew/PDF/hyperelliptics.pdf}.
\bibitem{Source27}
\emph{A defining polynomial in twenty-seven variables:
an existence theorem with a uniform degree bound},
companion paper, \nolinkurl{papers/few_variables27}, 9 November 2026.
\bibitem{Dependencies}
\emph{Arithmetic dependencies of the construction with 29 variables},
companion paper, \nolinkurl{papers/arithmetic_dependencies}.
It pins the source statements and Sun's prime-pattern theorem.
\bibitem{Baseline}
\emph{An explicit polynomial in 29 variables over the rationals},
companion paper, \nolinkurl{few_variables_20261108_v8.tex},
local implications and simultaneous-completeness sections.
\bibitem{Trace}
\emph{Aggregate trace heights of elliptic multiples}, companion
paper, \nolinkurl{papers/trace_heights}, applying J.~von Buhren,
\emph{J. Th\'eor. Nombres Bordeaux} \textbf{29} (2017), 307--320,
Theorem~1.1, \url{https://www.numdam.org/item/10.5802/jtnb.981.pdf}.
\bibitem{Logs}
\emph{Uniform formal-logarithm estimates}, companion paper,
\nolinkurl{papers/elliptic_logarithms}.
\bibitem{Active}
\emph{Separating the Nonzero Part of a Finite Algebra},
companion paper, \nolinkurl{active_finite_algebra_20261109_v15.tex}.
\end{thebibliography}
\end{document}
